% Supplementary Material — Paper 1 (English). Built before manuscript.tex (cross-references).
\documentclass[11pt]{article}
\newcommand{\DocLang}{en}
\input{papers/shared/preamble.tex}
\input{papers/shared/authors.tex}
\usepackage{pdflscape}
\hypersetup{pdftitle={Supplementary Material: Post-marketing safety of VA-MENGOC-BC in Cuba}}

\begin{document}
\SupplementaryNumbering
\SyntheticNotice
\begin{center}
{\Large\bfseries Supplementary Material\par}
\medskip
{\large Post-marketing safety of the VA-MENGOC-BC meningococcal B and~C vaccine in Cuba,
\VStudyFirstYear--\VStudyLastYear: a nationwide analysis of spontaneous reports with
disproportionality, latent class and machine-learning methods\par}
\medskip
{\small \AuthorOne{} and \AuthorTwo\par}
\end{center}

\tableofcontents

\section{Annual reporting rates}\label{sec:S-rates}
\begin{table}[htbp]
\centering\small
\caption{Annual AEFI reports involving VA-MENGOC-BC, doses administered and reporting rates per
\num{100000} doses with exact Poisson 95\% confidence intervals. The ratio compares each rate with
the programme's expected rate of \POExpectedRate{} per \num{100000} doses for all childhood
vaccines. Years without a dose denominator show a dash.}
\label{tab:S-rates}
\POTableRates
\end{table}

\section{Latent class model selection}\label{sec:S-lca}
\begin{table}[htbp]
\centering\small
\caption{Latent class models with $K=\POLCAKMin,\dots,\POLCAKMax$ classes fitted to the event indicators with a
prevalence of at least 0.5\%. The model with the lowest BIC was selected; ICL-BIC penalises
poorly separated classes and the entropy $R^2$ measures classification certainty (1 = perfect).}
\label{tab:S-lca}
\POTableLCASelection
\end{table}

\section{Correlates of hospitalisation}\label{sec:S-logistic}
\begin{table}[htbp]
\centering\small
\caption{Adjusted odds ratios of hospitalisation from the multivariable Firth logistic regression
fitted to all reports (Wald 95\% confidence intervals). Region contrasts are relative to the
western region; event and vaccine terms are present-versus-absent indicators. Terms recorded in
fewer than \VMinCell{} reports are suppressed.}
\label{tab:S-logistic}
\POTableLogistic
\end{table}

\section{Disproportionality under each comparator design}\label{sec:S-designs}
All tables follow the notation of the main text: $a$, reports with VA-MENGOC-BC and the event;
$E$, expected $a$ under independence; IC\textsubscript{025}, lower 95\% credibility bound of the
information component; EB\textsubscript{05}, lower 5th percentile of the empirical Bayes geometric
mean (primary design only, because the gamma Poisson shrinker prior is estimated on the whole
database); Methods, number of the four methods meeting their thresholds; Signal, prespecified
criterion (IC\textsubscript{025}${}>0$ with $a\geq\PODispMinReports$). Counts below \VMinCell{} are suppressed
together with the statistics derived from them.

\begin{table}[htbp]
\centering\small
\caption{Primary design: VA-MENGOC-BC reports against all other reports (all event types).}
\label{tab:S-disp-primary}
\POTableDispPrimaryAllOtherVaccines
\end{table}

\begin{table}[htbp]
\centering\small
\caption{Infant active-comparator design: reports in children younger than \POInfantMonths{}
months.}
\label{tab:S-disp-infant}
\POTableDispInfantActiveComparator
\end{table}

\begin{table}[htbp]
\centering\small
\caption{Single-vaccine design: VA-MENGOC-BC given alone against other vaccines given alone.}
\label{tab:S-disp-single}
\POTableDispExcludingCoadministration
\end{table}

\begin{table}[htbp]
\centering\small
\caption{Masking design: comparator group without the reports that mention the pentavalent
vaccine (PENTA-L).}
\label{tab:S-disp-masking}
\POTableDispExcludingPentavalentMasking
\end{table}

\section{Temporal stability of the information component}\label{sec:S-cumic}
\begin{table}[htbp]
\centering\small
\caption{Cumulative information component (IC) with 95\% credibility interval, recomputed with
the reports accumulated up to the end of each year, for the four most frequently reported events
and for every event meeting the primary signal criterion.}
\label{tab:S-cumic}
\POTableCumulativeIC
\end{table}

\section{Data quality}\label{sec:S-dq}
Data quality was assessed with the harmonised framework of Kahn et al.\ (conformance,
completeness and plausibility) on the harmonised reports before deduplication.

\begin{table}[htbp]
\centering\small
\caption{Completeness: percentage of reports with each key field recorded, by year of the source
file.}
\label{tab:S-completeness}
\POTableCompleteness
\end{table}

\begin{table}[htbp]
\centering\small
\caption{Plausibility checks over all years. When some yearly counts were suppressed, the number
flagged is shown as a lower bound and the percentage is omitted.}
\label{tab:S-plausibility}
\POTablePlausibility
\end{table}

\section{Meningococcal disease in Cuba, \POITSFirstYear--\POITSLastYear}\label{sec:S-its}
\begin{figure}[htbp]
\centering
\includegraphics{p1_incidence}
\caption{Annual national cases of meningococcal disease (logarithmic scale) and fit of the
segmented negative-binomial model. The shaded band marks the \POITSTransition{} transition (efficacy trial,
mass vaccination and inclusion in the national programme).}
\label{fig:S-its}
\end{figure}

\end{document}
